\problema{Coin Change}{322}{src/07-dp/322-coin-change.cpp}{M}

Dado un arreglo \texttt{coins} con las denominaciones de monedas disponibles
(cada denominación se puede usar tantas veces como se quiera) y un entero
\texttt{amount}, calcula el número mínimo de monedas necesarias para juntar
exactamente \texttt{amount}. Si no existe ninguna combinación que llegue
exacto, devuelve \texttt{-1}.

\paragraph{La idea, con un ejemplo de la vida real.} Piensa en dar cambio en
una caja registradora: quieres usar el menor número posible de monedas para
completar una cantidad exacta. En vez de probar a fuerza bruta todas las
combinaciones posibles de monedas (que crecen exponencialmente), resolvemos
primero los montos pequeños y reutilizamos esas respuestas para construir los
montos grandes. Esto es \textbf{programación dinámica} (\emph{dynamic
programming}, DP): guardamos la respuesta óptima de cada subproblema (cada
monto menor) para no recalcularla.

\begin{center}
\ttfamily
\begin{tabular}{l}
coins = [1, 2, 5], amount = 11\\
11 = 5 + 5 + 1 (3 monedas, es el mínimo)\\
\end{tabular}
\end{center}

\paragraph{Cómo lo resuelve el código, paso a paso.}
\begin{enumerate}
  \item Creamos una tabla \texttt{dp} de tamaño \texttt{amount + 1}, donde
        \texttt{dp[m]} representa el mínimo número de monedas para llegar
        exacto al monto \texttt{m}. Inicializamos todo en un valor
        ``infinito'' (representa ``todavía no sabemos cómo llegar a este
        monto''), salvo \texttt{dp[0] = 0} (con cero monedas llegas al monto
        cero).
  \item Para cada monto \texttt{m} de \texttt{1} a \texttt{amount},
        probamos cada moneda \texttt{c} del arreglo \texttt{coins}.
  \item Si \texttt{c} cabe en \texttt{m} (es decir \texttt{c <= m}) y ya
        sabemos cómo llegar a \texttt{m - c} (su valor en \texttt{dp} no es
        infinito), entonces una forma de llegar a \texttt{m} es usar esa
        solución y agregar una moneda \texttt{c} encima:
        \texttt{dp[m] = min(dp[m], dp[m - c] + 1)}.
  \item Al terminar, si \texttt{dp[amount]} sigue siendo infinito, no existe
        ninguna combinación exacta y devolvemos \texttt{-1}; si no,
        devolvemos \texttt{dp[amount]}.
\end{enumerate}

\lstinputlisting{src/07-dp/322-coin-change.cpp}

\complejidad{$O(\text{amount} \times k)$, con $k$ el número de denominaciones}{$O(\text{amount})$}

\paragraph{¿Qué significa esa complejidad?} Para cada uno de los
\texttt{amount} montos posibles probamos las $k$ denominaciones disponibles,
así que el trabajo total es proporcional al producto de ambos. El espacio es
lineal en \texttt{amount} porque solo guardamos una tabla de una dimensión
(no hace falta recordar qué moneda se usó en cada paso para resolver este
problema, solo el conteo mínimo).

\paragraph{Consejos para la entrevista (Amazon).} Es tentador proponer una
solución \textbf{greedy} (usar siempre la moneda más grande que quepa), pero
esa estrategia falla con denominaciones como \texttt{\{1, 3, 4\}} y
\texttt{amount = 6}: greedy elegiría \texttt{4 + 1 + 1} (tres monedas)
cuando la respuesta óptima es \texttt{3 + 3} (dos monedas). Vale la pena
mencionar ese contraejemplo para justificar por qué se necesita DP. También
conviene aclarar el caso límite \texttt{amount = 0} (la respuesta es
\texttt{0}, sin usar ninguna moneda) y el caso en que ninguna combinación
llega exacto (por ejemplo, monedas que solo suman números pares y un monto
impar).
